

\section{\textbf{研究背景与问题提出}}

科学研究是推动社会进步和技术创新的核心动力，而科研资助制度则是决定科研资源配置效率的关键。当前，国内外主流科研资助体系通常依赖同行评审（peer review）下的竞争性排序：评审专家基于对申请书科学价值、创新性与可行性的综合评估对项目进行排序，资助机构再依据排序结果确定最终资助对象。然而，随着科研竞争的加剧和申请数量的激增，该机制面临日益严峻的理论困境与实践挑战。

% \textbf{评审者面临的困境} \quad \textul{当申请数量庞大且整体质量趋同时，评审专家被迫在差异极小的申请之间做出细粒度区分，而这种区分往往超出了人类判断的一致性边界。}\cite{kahnemanNoiseFlawHuman2021} 系统论证了专家判断中噪声的普遍性，指出评审者间差异是评分异质性的主要来源。\cite{galloInfluencePeerReviewer2016}对AIBS (American Institute of Biological Sciences)的资助评审数据分析表明，评审分数中约 23\% 的异质性可归因于申请质量，剩余约 77\% 来自评审者间差异和噪声。\cite{pierLowAgreementReviewers2018}通过让多组评审专家独立评估同一批NIH（National Institutes of Health）的申请，发现不同评审组之间的评分一致性极低，进一步证实了评审噪声问题的严重性。\cite{schweigerCostsCompetitionDistributing2024}指出同行评审虽能区分“优秀”与“不合格”的申请，但在顶尖申请之间的甄别能力极为有限。评审噪声与一致性的不足削弱了精细排序的可信度，并提高了排序的边际成本。

% 更为关键的是，排序规则改变了评审者的决策框架与心理权重。研究表明，\textul{在难以精确比较项目优劣的情况下，评审判断更容易依赖申请人声誉、研究路径成熟度或项目成功确定性等并非直接反映科研潜力的信号，从而表现出偏见和风险规避倾向}\citep{fangResearchFundingCase2016,franzoniUncertaintyRisktakingScience2023}。当资助率较低、预算稀缺时，评审者面临一种高责任的决策环境：“资助了一个最终失败的项目”往往被视为可归责的显性错误，而“错过一个潜在突破性项目”则是不可观测的机会成本。在这种非对称的损失结构下，前景理论\citep{Tversky1992}所强调的损失厌恶与风险规避更容易被激活——评审者倾向于将“不确定性”编码为“潜在损失”，从而系统性低估高风险高回报项目的价值，偏好可验证、可预期的“安全型”路径。

\textbf{评审专家面临的困境} \quad \textul{当申请数量庞大且整体质量趋同时，评审专家被迫在差异极小的申请之间做出细粒度区分，而这种区分往往超出了人类判断的一致性边界。}研究表明，评审评分异质性的主要来源并非申请质量差异，而是评审者间的系统性分歧与随机噪声\citep{kahnemanNoiseFlawHuman2021, galloInfluencePeerReviewer2016, pierLowAgreementReviewers2018}。评审噪声削弱了排序的可信度，且排序越精细，为克服噪声所需付出的边际成本越高。更关键的是，排序规则改表了评审者的决策框架与心理权重。\textul{当项目优劣难以精确比较时，评审判断更易依赖申请人声誉、路径成熟度或成功确定性等信号，表现出偏见与风险规避倾向} \citep{fangResearchFundingCase2016, franzoniUncertaintyRisktakingScience2023}。当资助率较低、预算稀缺时，“资助一个最终失败的项目”是可归责的显性错误，而“错过一个潜在突破”则是不可观测的机会成本。这种非对称损失结构激活了前景理论\citep{Tversky1992}所强调的损失厌恶与风险规避——评审者倾向于将不确定性编码为潜在损失，系统性低估高风险高回报项目的价值。

% \textbf{申请人面临的困境} \quad 上述评审环境进一步塑造了申请人的策略行为。申请者作为理性的策略响应者，在排序竞争框架下可能产生两类典型的策略扭曲。\textul{其一，文书投入的军备竞赛。}\cite{holmstromMoralHazardObservability1979}的道德风险理论指出，当产出难以精确测度时，强激励可能诱导代理人将努力投入于可观测但非生产性的活动。在科研资助情境下，这意味着研究者可能将过多精力投入于申请书“包装”而非实际研究。\cite{schweigerCostsCompetitionDistributing2024}估计，撰写一份新申请通常需要25—50个工作日，而平均资助率仅为10\%—25\%，这意味着每资助一个项目需要消耗100—500人日的申请工作量。这种“军备竞赛”式的文书投入不仅造成巨大的经济成本，更挤占了实际研究时间，还可能导致研究人员心理健康问题、工作生活失衡以及不当研究行为。\textul{其二，风险规避与选题保守化。}  研究表明，低资助率下研究者会理性地放弃高风险高回报的研究选题\citep{grossContestModelsHighlight2019}，高度竞争的资助环境显著抑制了创新方向的选择 \citep{wangFundingModelCreativity2018}，而强排序竞争甚至可能导致高质量申请者的自我排除 \citep{addaGrantmakingGradingCurve2024}。

\textbf{申请者面临的困境} \quad 评审端的噪声与偏见进一步塑造了申请者的策略行为。在竞争性排序框架下，申请者可能产生两类典型的行为扭曲。\textul{其一，文书投入的军备竞赛。}当产出难以精确测度时，强激励可能诱导代理人将努力投入于可观测但非生产性的活动\citep{holmstromMoralHazardObservability1979}。科研资助情境下，研究者可能将过多精力投入于申请书“包装”而非实际研究，而在低资助率下每资助一个项目所消耗的申请工作量极为可观\citep{schweigerCostsCompetitionDistributing2024}，由此挤占研究时间并引发心理健康与不当研究行为等连带问题。\textul{其二，风险规避与选题保守化。}在低资助率与强排序竞争的环境下，研究者倾向于理性地回避高风险高回报的选题，转而选择更安全、可预期的研究路径，甚至可能导致高质量申请者的自我排除\citep{wangFundingModelCreativity2018, grossContestModelsHighlight2019, addaGrantmakingGradingCurve2024}。

\textbf{系统层面的恶性循环} \quad 评审者偏差与申请者策略相互强化，形成一个自我加强的负向循环：\textul{评审更保守$\to$申请更保守$\to$更多文书打磨、提案分布更趋同$\to$评审更依赖“成熟信号”与“安全路径”$\to$进一步压制高风险创新。该循环不仅解释了“安全型研究”激励的强化，也揭示了高强度竞争环境中文书投入膨胀与研究方向趋同的深层机制}。更为根本的问题是：\textbf{真正具有突破性的研究往往具有高不确定性，而这种研究往往在排序竞争中被系统性排斥。}\cite{whitleyImpactChangingFunding2018}通过跨国比较研究发现，竞争性项目资助占比的上升往往伴随资助框架的标准化，从而使偏离主流范式的研究越来越难以获得支持。\textul{这意味着，除效率损失之外，精细排序还可能从根本上抑制科学体系的创新性和多样性。}


% 在此背景下，\textbf{“合格筛选+随机抽签”（modified/partial lottery，以下称为“部分随机”）}\footnote{国内已有学者将类似机制概括为“底线评审”（刘欣 等, 2021; 鲍沁星 等, 2023），侧重评审端从“择优”向“择差”的转变。本项目采用“部分随机”，旨在同时涵盖评审端的简化与分配端随机性的引入，并与国际文献中 modified/partial lottery 的通行术语保持一致。}正在成为改革选项。其核心设计是：\textul{评审专家仅筛选出达到质量标准的合格项目，而不再进行精细化排序，最终资助对象通过随机方式从合格项目中产生。}这一设计试图在质量控制与效率之间取得平衡——保留同行评审的质量把关功能，同时避免在噪声中进行过度精细排序的效率损失。该机制具有多重潜在优势：降低边际排序所需的认知成本；抑制申请文书的军备竞赛，将努力由“包装”转向“实质科研”；缓解马太效应，为尚未崭露头角的研究者提供更多机会；并可能增强风险承担激励，从而促进更具创新性与突破性的研究。

% 在此背景下，\textbf{"合格筛选+随机抽签"（partial lottery，以下称为"部分随机"）}\footnote{国内已有学者将类似机制概括为"底线评审"（刘欣 等, 2021; 鲍沁星 等, 2023），侧重评审端从"择优"向"择差"的转变。本项目采用"部分随机"，旨在同时涵盖评审端的简化与分配端随机性的引入，并与国际文献中partial lottery的通行术语保持一致。}正在成为改革选项。其核心设计是：\textul{评审专家仅对申请进行合格/不合格的二元判定（binary screening），筛除不达标申请，最终资助对象通过从合格项目中以随机抽签产生方式产生}\citep{felicianiFundingLotteriesResearch2024}。该设计保留了同行评审的质量把关功能，同时避免在噪声主导区间进行无效排序，并可能带来多重收益：降低评审认知负荷、抑制文书军备竞赛、缓解马太效应、增强风险承担激励等。

在此背景下，\textbf{“合格筛选+随机抽签”（partial lottery，以下称为“部分随机”）}\footnote{国内已有学者将类似机制概括为“底线评审”\citep{BaochenzhouHRC2023, liu2021peerreview}，侧重评审端从“择优”向“择差”的转变。本项目采用“部分随机”，旨在同时涵盖评审端的简化与分配端随机性的引入，并与国际文献中partial lottery的术语保持一致。}正逐渐成为改革选项。其核心设计是：\textul{评审专家仅对申请作出合格/不合格的二元判定，据此剔除不达标申请，最终资助对象通过随机抽签的方式从合格项目中产生。}该设计在保留同行评审质量把关功能的同时，避免在噪声主导区间进行低效排序，并可能带来多方面收益，包括降低评审的认知负荷、抑制文书军备竞赛、缓解马太效应以及增强风险承担激励等。

% 然而，该机制对评审判断和申请人行为的影响仍缺乏系统研究。\textul{当评审专家知晓最终资助将通过随机方式产生时，其筛选行为是否会发生变化？申请人是否会因竞争框架的转换而调整选题风险与文书投入策略？机制转换最终能否改善资助组合的创新多样性？}这些问题直指“规则变化$\to$心理机制$\to$行为改变”的因果链条，却缺少直接的实验识别与可复制证据。因此，亟需开展以行为机制为核心的系统检验，为改革提供可操作的证据基础。

然而，关于部分随机机制如何影响评审者判断、申请人策略以及最终科研产出，现仍缺乏系统性的证据。\textul{当评审专家知晓最终资助将以随机方式决定时，其筛选标准与评判行为是否会相应调整？在竞争框架发生转换的情形下，申请人是否会改变选题风险取向与文书投入策略？这一机制转换最终能否提升资助项目的创新性与多样性？}上述问题缺少直接的实验识别与可复制的实证证据，因此亟需围绕相关行为机制开展进一步研究。

\section{\textbf{国内外研究现状及发展动态}}

\subsection{\textbf{同行评审中的评审者偏差}}

大量研究从不同角度揭示了同行评审中的系统性偏差：对新颖性与高风险项目的低估、对负面信息的过度反应、对申请者声誉与既往成就的依赖，以及由此引发的马太效应与累积优势。这些研究共同表明：在不确定性较高、信息不完备的环境中，评审者并非总能做出“中性、稳定、可重复”的质量判断。

\textbf{评审噪声} \quad 竞争性排序机制要求评审者完成的精细区分，在统计意义上缺乏可靠性。 \cite{kahnemanNoiseFlawHuman2021} 系统论证了专家判断中噪声的普遍性，指出评审者间差异是评分异质性的主要来源。 \cite{galloInfluencePeerReviewer2016}对AIBS (American Institute of Biological Sciences)的资助评审数据分析表明，评审分数中约 23\% 的异质性可归因于申请质量，剩余约 77\% 来自评审者间差异和噪声。\cite{pierLowAgreementReviewers2018}通过让多组评审专家独立评估同一批NIH（National Institutes of Health）的申请，发现不同评审组之间的评分一致性极低，进一步证实了评审噪声问题的严重性。\cite{schweigerCostsCompetitionDistributing2024}指出同行评审虽能区分“优秀”与“不合格”的申请，但在顶尖申请之间的甄别能力极为有限。\cite{liBigNamesBig2015}的研究虽然发现评审分数与未来产出之间存在正相关，但同时指出评审者难以有效预测哪些项目将产生突破性成果。

\textbf{保守倾向与损失厌恶} \quad 评审者往往表现出系统性的保守倾向，偏好可行性高、风险可控的申请。\cite{boudreauLookingLookingKnowledge2016}利用NIH创新研究项目的评审数据，发现提案与评审者知识距离越远、创新性越高，获得高分的概率越低，即使这些提案的预期科学价值更高。\cite{laneConservatismGetsFunded2022}设计了一项田野实验，通过在真实评审流程中随机操控申请书的负面信息暴露程度，证明评审者对负面信息赋予了过大的权重，呈现出系统性的保守偏差。\cite{franzoniUncertaintyRisktakingScience2023}从概念层面区分了科学研究中的“风险”（可概率化的不确定性）与“模糊性”（不可概率化的根本不确定性），指出当前评审体系对可行性的过度关注，实质上是将模糊性误作风险来处理，系统性低估了那些具有突破潜力但路径不确定的研究。

\textbf{光环效应与累积优势} \quad 评审者倾向于依据申请者的既往声誉而非项目本身质量进行判断，形成马太效应。 \cite{bolthijsMatthewEffectScience2018}对Netherlands Organization of Scientific Research 数据的分析显示，知名科学家的申请获得高分的概率显著高于同等质量的不知名科学家的申请。\cite{qiuMatthewEffectResearch2023}对NIH资助动态的研究表明，早期获得资助的研究者在后续竞争中具有显著优势，这种累积效应在职业生涯早期尤为明显。然而，目前尚无明确证据表明，这些累积优势最终会转化为相应的科研产出。\cite{mongeonConcentrationResearchFunding2016}的研究表明，获得最多资助的精英研究人员在产出和影响力方面并无突出优势。\cite{TianJiYuDIDMoXingDeKeJiZhengCeChuangXinNengLiZiZhuXiaoYingShiZhengYanJiuYiJieQingJiJinDiQiuKeXueXiangMuWeiLi2018}基于DID模型对杰青基金地球科学项目的评估发现，人才类资助对受资助者的创新能力具有显著正向效应，但这一效应主要惠及已获得资助的“赢家”，客观上加剧了资源集中。\cite{aagaardConcentrationDispersalResearch2020}通过系统性文献综述发现，科研资助的集中化趋势并未带来相应的产出效益提升，反而可能使边际收益递减。\cite{gubaEvaluatingGrantProposals2023}对俄罗斯科研资助数据的分析则从另一个角度揭示了这一问题：将基于发表指标的量化筛选引入评审流程后，虽有助于提高效率，但并未显著改善资助与未来研究产出之间的关联。

\subsection{\textbf{申请者对资助规则的策略反应}}

在申请者视角下，科研资助制度可被视为一类具有明确激励结构的竞赛（contest）。在申请成本为正且获资助概率不确定的情况下，研究者会依据评审标准与预期资助概率进行理性权衡。既有研究表明，资助规则不仅决定资源配置结果，还会通过影响进入决策、时间投入结构以及研究选题取向。


\textbf{自我选择} \quad \cite{addaGrantmakingGradingCurve2024}构建了科研资助中非市场资源配置的理论模型，并利用European Research Council评审规则变动作为自然实验加以检验。研究发现，当评审噪声上升时，申请数量反而增加；而当评审精度较高时，高质量申请者可能因预期竞争强度上升、净收益降低而自我排除。这一反直觉的结果源于申请者的策略选择：排序的不确定性越高，部分申请者反而更倾向于进入以期通过“运气”获得资助。

% \textbf{“包装”和实质科研的权衡} \quad \cite{schweigerCostsCompetitionDistributing2024}估计，撰写一份新申请通常需要25—50个工作日，而平均资助率仅为10\%—25\%，这意味着每资助一个项目需要消耗100—500人日的申请工作量。由此可见，在低资助率制度下，大量文书准备已构成显著的效率损失。在此基础上，\cite{myersTradeoffsCompetitionGrant2024}对科学家时间使用的调查研究揭示了科学家在筹资上花费的时间与其研究时间之间并非简单的替代关系—：一定程度的申请准备工作本身具有科学价值，但当竞争过度时，边际投入越来越多地转向非生产性的“包装”活动；并且当评审者难以准确评估项目质量时，申请书的“包装质量”可能与“研究质量”产生分离，导致“劣币驱逐良币”的逆向选择。

\textbf{“包装”与实质科研的权衡} \quad \cite{schweigerCostsCompetitionDistributing2024}估计，撰写一份新申请通常需要25—50个工作日，而平均资助率仅为10\%—25\%，即每资助一个项目需消耗100—500人日的申请工作量。然而，申请准备与研究时间并非简单的替代关系。\cite{myersTradeoffsCompetitionGrant2024}的调查研究表明，适度的申请准备本身具有科学价值，但当竞争过度时，边际投入越来越多地转向非生产性的“包装”活动。当评审者难以准确评估项目质量时，“包装质量”可能与“研究质量”产生分离，导致“劣币驱逐良币”的逆向选择。

% \textbf{自我选择} \quad \cite{addaGrantmakingGradingCurve2024} 构建了一个关于非市场资源配置的理论模型，用以解释科研资助中的资金分配机制。该模型表明，当某一领域的评审噪声上升时，该领域的申请数量反而会增加，而其他领域的申请数量相应减少；在极端情形下，当评审完全无噪声时，甚至可能出现高质量申请者的自我排除。作者利用European Research Council 评审规则的变动作为自然实验，实证检验并支持了上述预测。这一看似反直觉的结果源于申请者的策略选择：由于申请本身存在成本，当评审精度较高时，低质量申请者虽被有效筛除，但高质量申请者也因预期竞争强度上升、成功概率下降、预期净收益降低而选择退出；相反，当评审噪声增加、排序结果的不确定性提高时，部分申请者反而更倾向于进入，以期通过“运气”获得资助机会。


% \textbf{高风险选题} \quad \cite{grossContestModelsHighlight2019}构建的竞赛模型表明在低资助率情形下研究者会理性地放弃高风险研究主题，并且当研究者在职业压力等非科学价值动机驱动下申请资助时，竞赛的总体效率可能为负，即在扣除申请成本后，资助反而可能抑制科学进步。\cite{wangFundingModelCreativity2018}对日本科研体系的实证研究发现，竞争性资助对创新方向的抑制效应在低学术地位群体（如青年和女性研究者）中尤为显著，而稳定拨款（block funding）则为这些群体提供了追求新颖想法的安全空间。\cite{franssenDrawbacksProjectFunding2018}通过荷兰高绩效研究团队的案例研究表明，项目制资助相比于奖励和奖金等灵活资助形式，对研究过程施加了更多约束，在认识论层面上不利于偏离既有范式的探索。\cite{AWenDingBoKuanYuXiangMuZiZhuDeZiZhuChaiYiXingYanJiuJiYu20092018RiBenHuaXueLingYuZiZhuChanChuDeShuJuFenXi2021}基于日本化学领域的比较研究表明，稳定拨款与项目资助在研究团队构成、研究内容取向和产出绩效方面均存在显著差异——稳定拨款下的科研活动更关注研究问题本身，而项目资助则更多体现了科学共同体对遴选标准的适应性响应，提示过度依赖竞争性项目资助可能扭曲研究导向。

\textbf{选题方向与创新激励} \quad \cite{grossContestModelsHighlight2019}的竞赛模型表明，低资助率下研究者会理性地放弃高风险选题，并且当研究者在职业压力等非科学价值动机驱动下申请资助时，赛的总体效率可能为负，即在扣除申请成本后，资助反而可能抑制科学进步。 \cite{wangFundingModelCreativity2018}对日本科研体系的实证研究发现，竞争性资助对创新方向的抑制在低学术地位群体（如青年和女性研究者）中尤为显著，而稳定拨款则为其提供了追求新颖想法的空间。\cite{franssenDrawbacksProjectFunding2018}的荷兰案例研究表明，“资助项目”而非“资助人”在认识论层面不利于偏离既有范式的探索。\cite{AWenDingBoKuanYuXiangMuZiZhuDeZiZhuChaiYiXingYanJiuJiYu20092018RiBenHuaXueLingYuZiZhuChanChuDeShuJuFenXi2021}基于日本化学领域的比较研究进一步发现，稳定拨款下的科研活动更关注研究问题本身，而项目制资助则更多体现了对遴选标准的适应性响应，提示过度依赖竞争性排序可能系统性扭曲研究导向。\cite{liMotivatingInnovationImpact2024}发现中国优青资助虽提升了受资助者的产出数量与引用影响，却降低了研究的颠覆性指数（disruption index）。\cite{yangUnveilingImpactDual2024}对NIH与NSF（National Science Foundation）资助项目的大规模分析亦发现类似的“创新双重性”：受资助研究在事前新颖性上显著高于未资助研究，但在事后颠覆性上并无明显优势。\cite{LiGaoFengXianGaoHuiBaoXiangMuZiZhuYouZhuYuKeYanChuangXinMaJiYuKeXueDaShuJuDeWenXianOuHeWangLuoFenXi2025}对NIH“新创新者项目”的研究表明，传统竞争性资助容易引发科学家规避风险、开展保守性研究。



\subsection{\textbf{部分随机机制研究}}

% 随机抽签作为一种资源配置机制，在公共政策领域有着悠久历史。近年来，越来越多的学者与政策制定者开始探索将部分随机引入科研资助体系。

% \textbf{理论进展} \quad \cite{fangResearchFundingCase2016}首次系统论证了部分随机的可行性。\cite{depeuterModifiedLotteryFormalizing2022}则从认识论角度进一步论证，当评审不确定性较高时，抽签体现了对不确定性的诚实承认而非对评审的否定。在机制设计层面，\cite{ben-porathOptimalAllocationCostly2014}研究了在存在验证成本（costly verification）时的最优分配机制，证明当验证成本较高时，最优机制采取“阈值+抽签”的形式，为部分随机机制提供了严格的理论基础。\cite{carnehlDesigningScientificGrants2025}的最新研究进一步证明，在特定条件下，部分随机能够实现比精细排序更高的社会福利。\cite{felicianiFundingLotteriesResearch2024}构建了资助抽签的扩展分类体系，通过蒙特卡洛模拟评估了不同抽签类型在认识正确性、无偏公平性和分配公平性方面的表现。 \cite{goldbergPrincipledApproachRandomized2025}则提出了MERIT（Maximin Efficient Randomized Interval Top-k）方法，基于区间估计的原则性随机化选择框架，为在评审不确定性条件下实施随机化提供了可操作的算法工具。

% \quad 完全随机抽签（pure lottery）主张对所有申请者进行随机分配，不进行任何质量筛选。\cite{vaesenHowMuchWould2017}指出，在高度竞争的资助环境忠，如果评审成本过高而预测效度有限，完全随机可能是一种成本效益更优的替代方案。然而，完全随机忽视了质量差异，可能导致资源配置的严重扭曲，因此在实践中鲜有采纳。部分随机则主张在初步筛选的基础上进行随机分配，一定程度上兼顾质量控制与效率。

随机抽签作为一种资源配置机制，在公共政策领域有着悠久历史。\cite{greenberg1998chance}最早主张在合格申请者中直接抽签以替代昂贵且缓慢的同行评审。\cite{brezis2007focal}在此基础上提出了focal randomization策略：接受评审者一致认可的项目，拒绝一致认为不达标的项目，对意见不一致的项目实施随机资助。 \cite{fangResearchFundingCase2016}首次系统论证了部分随机的可行性， \cite{depeuterModifiedLotteryFormalizing2022}从认识论角度进一步指出，抽签体现了对评审不确定性的诚实承认而非对评审的否定。在机制设计层面， \cite{ben-porathOptimalAllocationCostly2014}研究了在存在验证成本（costly verification）时的最优分配机制，证明当验证成本较高时，最优机制采取“阈值+抽签”的形式，为部分随机机制提供了严格的理论基础；\cite{carnehlDesigningScientificGrants2025}进一步证明在特定条件下部分随机能实现比精细排序更高的社会福利。 \cite{felicianiFundingLotteriesResearch2024}构建了资助抽签的扩展分类体系，通过蒙特卡洛模拟评估了不同抽签类型在认识正确性与公平性方面的表现。\cite{goldbergPrincipledApproachRandomized2025}则提出了基于区间估计的MERIT（Maximin Efficient Randomized Interval Top-k）方法，为评审不确定性条件下实施随机化提供了可操作的算法工具。

% 新西兰健康研究委员会Explorer Grants计划（2013年启动，评审仅判断合格/不合格，合格者随机抽签资助）、德国大众基金会Experiment!计划（2017年启动，合格池中一半由评审团选择、一半由抽签决定）、瑞士国家科学基金会Postdoc.Mobility计划（2018年试行，2022年引入基于贝叶斯排名的统一评估程序）、以及奥地利科学基金1000 Ideas计划（2019年启动，评审判断与随机抽取混合设计）。

\textbf{国际实践探索} \quad 近年来，多个科研资助机构已将部分随机纳入正式流程。Health Research Council of New Zealand自2013年起在其Explorer Grants计划中仅要求评审判断申请是否合格，所有合格申请通过软件随机数抽签资助；Volkswagen Foundation于2017年在Experiment!计划中引入质量审查后的实体抽签环节，并设置wild-card机制允许评审团在抽签之外单独决定资助个别项目；Swiss National Science Foundation于2018年在Postdoc.Mobility计划中试行对评审未能区分优劣的合格申请采用抽签分配；Austrian Science Fund自2019年起在1000 Ideas Programme中采用混合设计，由评审直接选出部分高质量项目，其余同等价值的合格项目通过抽签决定。\footnote{关于上述四项计划的详细机制描述与比较，参见 \cite{bendiscioli2022experimental}。}这些实践在机制细节上各有差异，但共享“合格筛选+随机分配”的核心逻辑。\cite{gigerenzerAlternativeModelsFunding2025}在对上述实践及更广泛替代资助模型的综合讨论中指出，引入随机性是应对评审者间一致性偏低问题的一项更为根本性的改革方案之一；但同时也强调，目前尚不存在单一最优的好奇心驱动研究（curiosity-driven research）资助机制，且对替代模型开展系统性实证评估的研究仍显不足。

\textbf{创新激励的复杂效应} \quad 部分随机对创新行为的影响并非单向度。一方面，随机性降低了资助结果的可预测性，可能削弱过度包装与风险规避的动机。\cite{mansoMotivatingInnovation2011}指出对早期失败的容忍是激励创新的关键机制，\cite{edererPayPerformanceDetrimental2013}在实验室环境中验证了这一预测，发现“容忍早期失败+奖励长期成功”的激励方案显著提升了探索行为。\cite{azoulayIncentivesCreativityEvidence2011}通过比较霍华德·休斯医学研究所（Howard Hughes Medical Institute, HHMI）研究员与NIH资助者进一步表明，容忍失败的资助模式能显著促进高影响力成果与探索性创新。另一方面，若申请者将结果理解为“主要由运气决定”，则可能削弱其投入高质量申请的激励，诱发逆向选择。\cite{charnessCreativityIncentives2019}的实验研究则指出，激励机制对创造力的影响取决于任务结构：竞赛激励在约束较强的封闭型任务中有助于提升创造力，而在开放型任务中效果不显著，同伴排名等非金钱激励反而更为有效。这些证据提示，部分随机对创新的影响可能取决于研究类型、任务结构与激励强度之间的匹配关系，有必要通过更精细的机制设计与实验识别加以检验。



% \textbf{创新激励的复杂效应} \quad 一方面，“随机”降低了资助结果的可预测性，从而可能削弱申请者进行过度包装与风险规避的动机，鼓励更加真实且具有探索性的研究提案。 \cite{mansoMotivatingInnovation2011}指出，对早期失败的容忍是激励创新的关键机制。\cite{edererPayPerformanceDetrimental2013}在实验室环境中验证了这一理论预测，发现“容忍早期失败 + 奖励长期成功”的激励方案显著提升了被试的探索行为。\cite{azoulayIncentivesCreativityEvidence2011}通过比较HHMI（Howard Hughes Medical Institute）研究员与 NIH 资助者的实证分析进一步表明，容忍早期失败的资助模式能够显著促进高影响力成果与探索性创新的产生。另一方面，若申请者将结果理解为“主要由运气决定、与质量关联较弱”，则可能削弱其投入高质量申请的激励，诱发逆向选择。\cite{liMotivatingInnovationImpact2024}对中国国家自然科学基金优秀青年科学基金项目的研究发现，该资助虽提升了受资助者的产出数量与引用影响，但却降低了研究的颠覆性指数（disruption index），表明人才类资助可能促使研究者采取更偏向利用而非探索的策略。\cite{yangUnveilingImpactDual2024}对 NIH 与 NSF（National Science Foundation）资助项目的大规模分析亦发现类似的“创新双重性”现象：受资助研究在事前新颖性上显著高于未资助研究，但在事后颠覆性上并无明显优势，说明资助制度在激励新颖组合的同时，未必能够转化为范式突破。\cite{LiGaoFengXianGaoHuiBaoXiangMuZiZhuYouZhuYuKeYanChuangXinMaJiYuKeXueDaShuJuDeWenXianOuHeWangLuoFenXi2025}的NIH的“新创新者项目”研究表明，传统竞争性项目资助容易引发科学家规避风险、开展保守性研究的问题。\cite{charnessCreativityIncentives2019}的实验研究则进一步指出，激励机制对创造力的影响取决于任务结构：在约束较强的“封闭型”任务中，竞赛激励有助于提升创造力产出；而在开放型任务中，竞赛激励并不显著，同伴排名等非金钱激励反而更为有效。这些证据提示，部分随机机制对创新行为的影响并非单向度，而可能取决于研究类型、任务结构与激励强度之间的匹配关系，因而有必要通过更精细的机制设计与实验识别加以检验。



尽管随机资助机制在理论与实践层面均取得进展，但现有研究存在三方面显著局限。第一，缺乏对部分随机机制行为效应的系统研究。现有关于改良抽签的讨论多为规范性分析或模拟研究，缺乏基于行为数据的实证检验。\cite{barnettImpactWinningFunding2024}的新西兰HRC（Health Research Council of New Zealand）随机对照试验虽发现随机资助与排序资助在产出效果上无显著差异，但样本量有限（88名研究者），且未能系统测量申请者的心理与行为反应。第二，理论模型与实验验证的结合不足。既有研究或侧重于理论建模\citep{ben-porathOptimalAllocationCostly2014,grossContestModelsHighlight2019,carnehlDesigningScientificGrants2025}，或侧重于事后数据的实证描述\citep{barnettImpactWinningFunding2024,liMotivatingInnovationImpact2024}，缺乏将理论预测与控制实验有机结合的研究设计。第三，对中国情境的适用性研究缺失。现有研究主要基于欧美发达国家的资助体系，对中国科研资助环境的特殊性——如高竞争率、强绩效考核、人才项目与项目资助并存等制度特征——关注不足。

\section{\textbf{中国情境与现实需求}}

我国科研资助体系规模大、竞争强、评价压力高，排序化评审与绩效考核相互强化，使青年科研人员尤其容易被卷入“多轮申请—高投入—低命中率”的重复博弈。\cite{DuWoGuoKeJiZhengCeXueYanJiuTaiShiJiGuoJiBiJiao2017}指出，中国科技政策研究面临从“科技政策”向“科技创新政策”转型的迫切需求，而资助机制改革正是这一转型的核心议题。

\textbf{高竞争率与“内卷化”} \quad 国家自然科学基金的申请量近年来持续上升。以面上项目为例，如图1所示，2020年申请数量才超过10万项，至2025年已接近20万项。与此同时，资助率与资助强度呈持续下降趋势：资助率由2020年的17.15\%降至2025年的11.56\%（其中医学科学部的资助率更是低至8.92\%）；资助强度则由57.5万元降至50万元以下。这一高度“内卷”的竞争环境在很大程度上塑造了科研人员的行为模式。
\begin{figure}[h!]
	\centering %图片居中
	\includegraphics[width=\textwidth]{figures/trend.jpg}
	\captionsetup{justification=centering} %图题居中
	\caption{2020—2025年国自科面上项目受理申请数、资助率与资助强度变化}
    \vspace{4mm}
    \scriptsize{数据来源：国自科基金委2020-2024年年度报告及2026年度项目指南}
\end{figure}

\textbf{强绩效考核与行政化评审} \quad 我国科研评价体系长期以来侧重量化产出与短期绩效，“唯论文、唯职称、唯学历、唯奖项”的倾向尤为明显。\cite{xuImpactsIncentivesInternational2021}基于中国人文社会科学领域的证据指出，以国际发表数量为核心的激励机制在科研文化中引致多重张力，具体表现为国际化与本土化、数量与质量以及学术诚信与功利主义之间的冲突。并且，“一刀切”式激励政策可能进一步强化管理主义取向并削弱研究者的学术自主性。此类“锦标赛制”评价与精细排序的资助配置机制相互强化，从而共同塑造科研人员的保守型激励。

\textbf{青年与女性科研人员的特殊困境} \quad 在现行体系下，青年科研人员面临“非升即走”的聘期考核压力与基金申请的累积优势效应的双重挤压。\cite{HuangWoGuoNuXingKeYanRenYuanFaZhanXianZhuangTiaoZhanJiZhengCeYanBian2018}和\cite{YuKeXueJiJinZhuLiNuXingKeYanRenYuanChengChangZhengCeChengXiaoYuZhanWang2023}的研究分别从发展现状和政策演变角度表明，女性科研人员在获得科研资助和奖励方面的代表性仍然不足，基金资助体系虽在持续完善，但性别差异的结构性根源尚未完全消除。\cite{LiGaoFengXianGaoHuiBaoXiangMuZiZhuYouZhuYuKeYanChuangXinMaJiYuKeXueDaShuJuDeWenXianOuHeWangLuoFenXi2025}的研究进一步发现，青年科学家和女性科学家能更充分地发挥高风险高回报项目对创新产出的促进效应，表明这些群体的创新潜力在现行竞争性资助体系中可能未被充分释放。

基金委近年来推进“负责任评审”、“分类评审”、“分阶段评审”等改革，强调减负增效与质量导向，为探索“减少排序压力、优化程序设计”的机制创新提供了政策窗口。然而，在中国情境下，关于“减少排序、引入随机性”的讨论仍面临三类现实阻力：其一，担忧质量与激励受损——若抽签结果与质量无关，是否会削弱申请者投入高质量申请的动机？其二，担忧程序不透明引发信任下滑——若抽签过程缺乏可验证的透明度，申请者是否会质疑存在暗箱操作？其三，担忧评审者在筛选阶段出现“标准松动”——若知晓最终资助将通过抽签决定，评审者是否会在质量把关环节降低阈值？这些阻力的背后，核心在于缺乏足以令人信服的微观行为证据。因此，有必要在中国可识别的制度与文化语境下，开展以行为机制为核心的系统检验，为改革提供可操作、可验证的证据基础。

\section{\textbf{拟解决的科学问题}}

本项目围绕"规则变化$\to$心理机制$\to$行为反应"的因果链条，拟解决以下三个科学问题。

\noindent\textbf{Q1：评审者行为——排序压力如何塑造评审偏差？机制转换能否缓解？}

% 在评审噪声不可避免的条件下，排序竞争是否会放大评审者的损失厌恶与认知负荷，从而诱发更强的保守倾向与对"安全信号"的依赖？若将评审任务从精细排序转为合格筛选，这些偏差能否得到缓解？其缓解的心理来源是减少了边际比较的认知成本，还是降低了策略性评分的动机？

\noindent\textbf{Q2：申请者行为——机制转换如何改变选题风险与努力配置？}

% 当资助规则从排序竞争转向合格筛选时，申请者的决策参照点是否由“击败对手”转向“达到阈值”，其损失厌恶与风险权重是否发生系统性变化？这种心理转变能否导致更少的包装性投入与更高的探索性选题比例？

\noindent\textbf{Q3：成果结构——机制转换能否改善资助组合的创新多样性？其最优设计参数的边界条件是什么？}

% 本问题在Q1和Q2的行为证据基础上，考察资助组合层面的结构性指标（创新多样性、风险型项目占比、总体效率），并进一步识别不同设计参数组合（随机比例、合格阈值、程序透明度与问责机制）对上述指标的影响，确定机制优化的边界条件。



% \textcolor{red}{=======}

% 免责声明：敬请大家仔细对比本模版与官方word转pdf后的差别，自行确定是否采用。
% 如果采用博后的分阶段评审，限制只能用word，那么本项目就寿终正寝了。个人认为，2026年系统仍然上传PDF，只要人眼无法分辨与官方区别，就可以。

% \textbf{只需要修改`sections`,`figures`,`ref.bib`，那么如果模版更新，只需要拷贝这三项过去直接编译，就是最新版。}

% \section{自定义设置方法}
% 【常规】自定义设置方法

% \textbf{只需要修改`sections/个性化设置.sty`文件。}
% \begin{lstlisting}[language=tex, basicstyle=\ttfamily\small, keywordstyle=\color{blue}, commentstyle=\color{gray}]
% %【1】标题颜色
% \definecolor{HighLight}{RGB}{0,0,0}
% %【2】正文字体设置
% %\global\MS@kaitrue\global\MS@songfalse%楷体正文
% \global\MS@kaifalse\global\MS@songtrue%宋体正文
% %【3】是否屏蔽正文
% \global\MS@offtrue\global\MS@onfalse%屏蔽正文
% %\global\MS@ontrue\global\MS@offfalse%不屏蔽正文
% \end{lstlisting}

% 【高阶】
% 需要修改`nsfc.sty`文件中的定义。**如非必要，不建议修改。**

% \section{编译方法}
% \begin{enumerate}
% 	\item 编译：XeLaTex->bibtex->XeLaTeX->XeLaTeX
% 	\item 排错：多看编译错误，多查询错误解决方法；编译警告，只要不影响PDF，就不用管。本模版多人使用，可以认为不存在编译错误。
% \end{enumerate}


% \section{编辑方法}
% %%%%%%%%
% %\newpage
% \vspace{-5pt}

% \begin{figure}[h!]
% 	\centering %图片居中
% 	\includegraphics[width=6cm]{figures/xiugai.png}
% 	\captionsetup{justification=centering} %图题居中
% 	\caption{项目文件夹结构}
% \end{figure}
% 只需要修改图中蓝色选中区域文件，即\verb|sections, figures, ref.bib|即可。其他剩余文件，可以直接采用模版替换，编译就是最新版。

% \subsubsection{章节}\label{subsubsec:t}

% \begin{itemize}
% 	\item \verb|\label{subsubsect:t}|
% 	引用章节\verb|\ref{subsubsec:t}|，生成为：\ref{subsubsec:t}。这个样式可能不是你想要的，那种情况下，就手敲吧！申请书不像论文，这种情况应该没几个。
% 	\item 章节标题的话，如果编号、字体啥的不喜欢，也可以手动敲。示例如下：
	
% 	\noindent\textcolor{blue}{\textbf{1.1 测试}}
	
% 	其中，\verb|\noindent|标识没有首行缩进。要控制缩进距离，可以用\verb|\hspace{10pt}|手动缩进。缺点就是，不能自动编号，不能自动引用。但是，申请书不像论文，这种情况应该没几个。
% \end{itemize}

% \subsubsection{字体}
% 中文\textbf{粗体}；\textbf{bold} font；
% 中文\textit{斜体}；\textit{italic} font；

% 克制使用以下标注（不用更好），防止专家眼花缭乱。\textul{添加下划线}；\textull{添加双下划线}；\textuw{添加下弯线}；\textud{添加下点线}。

% 全文改宋体，可以修改nsfc.sty的MS部分字体。

% 可选的就是\verb|\zhkai,\enkai,\zhsong,\ensong|。

% \subsubsection{文献}
% 普通引用\cite{test}；上标引用\citess{test}；多篇文章\citess{test,test2,test3}。

% 有注音的英文：\cite{test}。\citep{test}

% 参考期刊\cite{test}；
% 参考图书\cite{test2}；
% 参考会议\cite{test5}；
% 参考链接\cite{test4}；
% 参考文件\cite{test6}。

% 对于中文参考文献，bib条目中需要有language = {zh}，参见\cite{test2}。

% 对于非数字形式/作者年份形式的参考文献引用，阅读nsfc.sty对应这一部分进行修改：
% \begin{lstlisting}[language=tex, basicstyle=\ttfamily\small, keywordstyle=\color{blue}, commentstyle=\color{gray}]
% \usepackage[square,numbers,sort&compress]{natbib}%数字引用[1]
% \bibliographystyle{bibs/gbt7714-numerical}
% %\usepackage[round,sort]{natbib}%作者年引用Author(2022)
% %\bibliographystyle{bibs/gbt7714-author-year}
% %bst定制方法：修改对应.bst文件的FUNCTION {load.config}部分设置
% \end{lstlisting}

% 要进一步定制bst的话，需要修改对应.bst文件的FUNCTION {load.config}部分设置。
% 例如，如果需要设置是否要文献类型标识，可以修改bibs/gbt7714-numerical.bst文件或bibs/gbt7714-author-year.bst文件的以下部分：
% \begin{lstlisting}[language=tex, basicstyle=\ttfamily\small, keywordstyle=\color{blue}, commentstyle=\color{gray}]
%   #0 'show.mark := %【可修改】如果不想要[J][M]这种文献类型标识，改为#0；想要文献类型标识改为#1
% \end{lstlisting}

% \subsubsection{列表}
% 无序列表\footnote{值得注意的是，不需要一定要用列表环境，用加粗、换行、缩进同样能达到效果。
% 	因为咱们的初衷，还是LaTeX在排版文献和公式上有优势，发挥这一个优势就行了，其他部分不需要强行套用。文本本身还是最重要、需要大家投入精力的部分。}的例子：
% \begin{itemize}[left= 50pt]
% 	\item[-] 第一条，第一条的内容可能很长长长长长长长长长长长长长长长长长长长长；
% 	\item[-] 第二条。
% \end{itemize}

% 有序列表的例子：
% \begin{enumerate}[left= 50pt]
% 	\item 第一条，第一条的内容可能很长长长长长长长长长长长长长长长长长长长长；
% 	\item 第二条。
% \end{enumerate}

% 两个带圈文字的实现方法：
% \textcircled{\raisebox{-0.8pt}{1}}
% \textcircled{\textbf{\small 1}}

% 注意，由于列表的缩进，不同使用者可能偏向并不一样。本模版用的enumitem包，阅读他的文档进行个性化，其文档在：https://www.ctan.org/pkg/enumitem


% \subsubsection{公式}

% 数学字体：
% \verb|\mathcal|：$\mathcal{XYZ},\mathcal{H}^T$，
% \verb|\mathbb|：$\mathbb{XYZ}$，
% \verb|\mathbf|：$\mathbf{XYZ}$，这个往往不加粗，需要用
% \verb|\symbf|：$\symbf{XYZ}$，
% \verb|\symbfit|：$\symbfit{XYZ}$

% \begin{itemize}
% 	\item 标量/普通变量：$x,t$
% 	\item 向量：$\symbfit{v}$；矩阵：$\symbfup{A}$
% \end{itemize}

% 编号公式如下：（这个圆括号的内容只是为了方便看公式与文字的行间距）
% \begin{align}
% \begin{cases}
% a\\b
% \end{cases}
% \end{align}
% 以及（这个圆括号的内容只是为了方便看公式与文字的行间距）
% \begin{equation}
% E = mC^2
% \end{equation}
% \begin{align}
% 	\max_{x \in C} g(x).
% \end{align}
% 不编号公式如下：（这个圆括号的内容只是为了方便看公式与文字的行间距）
% \begin{equation*}
% 	\begin{cases}
% 		a\\b
% 	\end{cases}
% \end{equation*}
% \verb|$$x=y$$|类型公式：（半行间距）
% $$
% 	\begin{cases}
% 		a\\b
% 	\end{cases}
% $$
% 公式的上下间距参见nsfc.sty中公式上下间距部分。




% \subsubsection{图}
% 图片的例子：
% \begin{figure}[h!]
% \centering %图片居中
% \includegraphics[width=2cm]{figures/xiugai.png}
% \captionsetup{justification=centering} %图题居中
% \caption{这是图题This is caption}
% \end{figure}

% 图题和表头若想取消加粗，去掉nsfc.sty中caption部分的\verb|\bfseries|即可。


% \subsubsection{表}
% 在表格内的第一行设置\verb|\zhkai\ensong\selectfont|，来选择字体。

% 其中\verb|\zhkai\zhsong\enkai\ensong|可以根据需要选择。
% \begin{table}[htbp]
% 	\zhkai\ensong\selectfont%设置表格字体
% 	\centering  % 显示位置为中间
% 	\caption{表格}  % 表格标题
% 	\label{table1}  % 用于索引表格的标签
% 	%字母的个数对应列数，|代表分割线
% 	% l代表左对齐，c代表居中，r代表右对齐
% 	\begin{tabular}{|c|c|c|c|}  
% 		\hline  % 表格的横线
% 		& & & \\[-6pt]  %可以避免文字偏上来调整文字与上边界的距离
% 		第一列&第二列&第三列&第四列 \\  % 表格中的内容，用&分开，\\表示下一行
% 		\hline
% 		& & & \\[-6pt]  %可以避免文字偏上 
% 		0.1&0.2&0.3&0.4 \\
% 		\hline
% 	\end{tabular}
% \end{table}

% \subsubsection{定义定理等}
% \begin{definition}
% 这是一个定义
% \end{definition}
% \begin{theorem}
% 	这是一个定理
% \end{theorem}

% \section{某页最后一段行距可能很窄？}
% 如果没有这个问题，就不用管这个事情。

% 行间距变化一般是在“多行蓝色模版”部分前后。因为蓝色模版文字在section里写的，latex把蓝色部分当作一个整体，可能硬要挤到这一页，而不是换新一页，导致会挤前一页的行间距，导致前一页行距异常。
% 针对这种情况，模版已经使用
% \begin{lstlisting}[language=tex, basicstyle=\ttfamily\small, keywordstyle=\color{blue}, commentstyle=\color{gray}]
% 	%自动段落的行间距微调
% 	\usepackage{setspace}
% 	\setstretch{1.6} % 22 bp / 14 pt = 1.571
% \end{lstlisting}
% 降低了这种情况发生的可能。
% 如果还有，就只好添加\verb|\newpage|把它newpage到后一页上，就行了。也可以考虑分段缓解，需要写的时候注意页面的分段和字数。

% \textcolor{red}{=======}

\begin{REF}
\subsection*{参考文献}
\vspace{-50pt}
\bibliography{ref}%参考文献
\end{REF}

% \newpage%自己判断是否需要
